\documentclass{beamer}

\usepackage{verbatim}
\usepackage{graphicx}
\usepackage{booktabs}
\usepackage[brazil]{babel}
\usepackage[utf8x]{inputenc}
\usepackage[portuguese, ruled, linesnumbered]{algorithm2e}
\usepackage{float}
\usetheme{Madrid}
\usepackage{tikz}
\usepackage{listings}
\usepackage{float}
\usepackage{adjustbox}
\usepackage{amsmath}
\usepackage{geometry}
\usepackage{colortbl}

\newcommand{\cinza}{\cellcolor[gray]{0.5}}
\definecolor{gray}{gray}{0.6}


\title[Aula 01 - Instalação]{Aula 01 - Instalação KiCad}

\author[CAECP]{Centro Acadêmico de Engenharia de Computação}

\institute[UTFPR] 
{
  Departamento de Informática - DAINF \\
  Universidade Tecnológica Federal do Paraná
}

\date[\today]{Mini-curso KiCad\\\today}

\begin{document}
\begin{frame}
  \titlepage
\end{frame}

\begin{frame}{Sumário}
  \tableofcontents
\end{frame}

\section{Introdução}
\begin{frame}{Introdução}
	\begin{columns}
		\begin{column}{0.5\textwidth}
			\begin{itemize}
				\item Software para elaboração de placas de circuito impresso, as PCBs.
				\item Ferramenta \textit{Open-Source}.
			\end{itemize}
			\begin{figure}
				\centering
				\includegraphics[width=1\textwidth]{Imagens/exemplo.png}
				\caption{Exemplo de um projeto \textit{open hardware}.}
			\end{figure}
		\end{column}
		\begin{column}{0.5\textwidth}
			\begin{figure}
				\centering
				\includegraphics[width=1.0\textwidth]{Imagens/KiCad-Logo.eps}
				\caption{Logo KiCad.}
			\end{figure}
		\end{column}
	\end{columns}
\end{frame}

\section{Instalação e Download}
\begin{frame}{Instalação e Download}
	\begin{block}{Windows 64 bits}
		\textit{Download} pelo \textit{link}:
		\begin{itemize}
			\item \url{http://downloads.kicad-pcb.org/windows/stable/kicad-5.0.0-x86\_64.exe}
		\end{itemize}
		Depois de baixado, fazer a instalação clássica.\\
		Não é preciso fazer instalação de biblioteca adicional.
	\end{block}
	\begin{block}{Ubuntu}
		Pelos comandos no terminal:
		\begin{itemize}
			\item \texttt{sudo add-apt-repository --yes ppa:js-reynaud/kicad-5}
			\item \texttt{sudo apt update}
		\item \texttt{sudo apt install kicad}
		\end{itemize}
		Não é preciso fazer a instalação de biblioteca adicional.
	\end{block}
\end{frame}

\begin{frame}{Instalação e Download}
	\begin{block}{Arch Linux}
		Pelo comando:
		\begin{itemize}
			\item \texttt{sudo pacman -S kicad}
		\end{itemize}
		Para instalar as bibliotecas adicionais (deixando a instalação completa):
		\begin{itemize}
			\item \texttt{sudo pacman -S kicad-library kicad-library-3d}
		\end{itemize}
	\end{block}
\end{frame}

\begin{frame}
  \titlepage
\end{frame}

\end{document}



